\documentclass[11pt,a4paper]{letter}
\usepackage[T1]{fontenc}
\usepackage[utf8]{inputenc}
\usepackage[margin=25mm]{geometry}
\usepackage{hyperref}
\signature{Vincenzo Barone}
\address{Consorzio Interuniversitario Nazionale per la Scienza e Tecnologia dei Materiali\\Italy}
\begin{document}
\sloppy
\begin{letter}{Editor-in-Chief\\International Journal of Quantum Chemistry}
\opening{Dear Editor,}

I am pleased to submit the manuscript entitled
\emph{“Semiempirical Orbital-Optimized DFT in MOUSE: Spin-Pure Excited-State
Surfaces with Analytic Gradients and Hessians”} for consideration as a Full
Paper in the \emph{International Journal of Quantum Chemistry}.

The manuscript addresses a practical limitation of excited-state quantum
chemistry: repeated calculations on large molecules require a state-specific
energy surface with reliable gradients and curvatures, while conventional
response methods are primarily designed for efficient vertical multi-root
spectra. We develop a DFT-informed semiempirical orbital-optimized model in
MOUSE that combines a one-time DFT reference with a spin-polarized GFN2
valence calculation, spin-pure singlet configurations, selected screened
double excitations, and complete analytic response through the nuclear
Hessian.

The principal results are: (i) a common-orbital spin-pure excited-state
surface with explicit state following and no fitted excitation-energy offset;
(ii) analytic gradients and Hessians including orbital stationarity, spin
recoupling, configuration-space response, the moving ground projector and
state-eigenvector response; (iii) a 26-state development assessment with
0.183 eV mean absolute error, including a marked improvement for the
$\pi\rightarrow\pi^*$ subset over a matched TDA control; and (iv) larger
carotene and chlorophyllide examples demonstrating configuration-space
tracking and convergence diagnostics. The manuscript is positioned relative
to simplified TD-DFT/sTDA, TD-DFTB, $\Delta$DFTB, TDDFT+TB and semiempirical
multireference alternatives, while clearly separating demonstrated results
from future independent-transferability tests.

This work is appropriate for IJQC because it presents a quantum-chemical
methodological development with an explicit variational formulation,
analytic nuclear response and an implementation intended for molecular
simulation and multilevel QM/QM$'$ applications. The accompanying Supporting
Information contains the full protocol, provenance, machine-readable tables,
derivative checks and additional convergence analyses.

The manuscript is original, has not been published in full, and is not under
consideration elsewhere. The author declares no conflict of interest. All
data and source code required to reproduce the reported results are available
in the MOUSE repository; large archived checkpoints are retained with
checksummed provenance and will be deposited in a permanent public archive
before publication.

Suggested reviewers are experts in excited-state electronic structure,
orbital-optimized DFT and semiempirical quantum chemistry. We would be happy
to provide specific names and contact details through the submission system.
We have no non-preferred reviewer requests.

Thank you for considering this manuscript. I look forward to your response.

\closing{Sincerely,}
\end{letter}
\end{document}
